% CISC 322/326 Assignment 2: Concrete Architecture report, Group WEZER.
% Compile with pdfLaTeX or XeLaTeX (Overleaf and Prism defaults work). Each section lives in sections/,
% numbered as in the group's work breakdown; write only in the section file you own.
%
% While \outlinetrue is set, the document prints the outline notes before the title page, an
% owner and target line under each heading, and a bulleted requirements list in every section.
% Write your prose below the list and delete the list once the section is done.
% \outlinefalse hides all of the outline material at once.
\documentclass[10pt,letterpaper]{article}

% ---------- outline switch ----------
\newif\ifoutline
\outlinetrue

% ---------- page and fonts ----------
\usepackage{iftex}
\ifPDFTeX
  \usepackage[T1]{fontenc}
  \usepackage[utf8]{inputenc}
  \usepackage{lmodern}
\fi
\usepackage[margin=1in]{geometry}
\linespread{0.96}
\usepackage{microtype}

% ---------- figures, tables, references ----------
\usepackage{graphicx}
\graphicspath{{figures/}}
\usepackage{booktabs}
\usepackage{tabularx}
\usepackage{array}
\usepackage[font=small,labelfont=bf]{caption}
\usepackage[section,above]{placeins}% keep each figure and table inside its own section
\usepackage{flafter}% a float never appears before the text that introduces it
\usepackage{float}
\usepackage{enumitem}
\usepackage{xurl}

% ---------- diagrams ----------
% UML diagrams are TikZ files in figures/. They were laid out with Helvetica-width text, so
% they are drawn in Helvetica (pdfLaTeX) or its clone TeX Gyre Heros (XeLaTeX).
\usepackage{tikz}
\usetikzlibrary{arrows.meta,backgrounds,fit,positioning}
\ifPDFTeX
  \newcommand{\diagramfont}{\fontfamily{phv}\selectfont}
\else
  \usepackage{fontspec}
  \newfontfamily\diagramfont{texgyreheros}[Extension=.otf, UprightFont=*-regular,
    BoldFont=*-bold, ItalicFont=*-italic, BoldItalicFont=*-bolditalic]
\fi
% \diagram{<file>.tex}: a TikZ diagram from figures/, scaled to the text width.
\newcommand{\diagram}[1]{\resizebox{\linewidth}{!}{\diagramfont\fontsize{11}{13}\selectfont\input{figures/#1}}}

% ---------- headings ----------
% One step above body text for sections, body size for subsections (as in KOMA-Script's
% headings=normal, memoir's \reparticle and llncs); emphasis comes from bold and space.
\usepackage{titlesec}
\titleformat*{\section}{\large\bfseries}
\titleformat*{\subsection}{\normalsize\bfseries}
\titleformat*{\subsubsection}{\normalsize\itshape}
\titlespacing*{\section}{0pt}{2.5ex plus 1ex minus .2ex}{1ex plus .2ex}
\titlespacing*{\subsection}{0pt}{2ex plus 1ex minus .2ex}{0.6ex plus .2ex}

% ---------- outline material ----------
% \sectionnote{...}: owner, page target and rubric criteria under a section heading.
% requirements environment: the bulleted list of what the section must contain.
% Both print only while \outlinetrue is set.
\usepackage{environ}
\newcommand{\sectionnote}[1]{\ifoutline\par\noindent{\small\itshape #1}\par\smallskip\fi}
\NewEnviron{requirements}{%
  \ifoutline
    \par\begingroup\small\raggedright
    \begin{itemize}[leftmargin=1.4em,itemsep=1pt,topsep=3pt,parsep=0pt]\BODY\end{itemize}%
    \endgroup\par
  \fi}

\usepackage[hidelinks]{hyperref}% load last

% ---------- report details ----------
\newcommand{\reporttitle}{WezTerm Concrete Architecture}
\newcommand{\reportcourse}{CISC 322/326 Software Architecture}
\newcommand{\reportassignment}{Assignment 2: Concrete Architecture}
\newcommand{\reportgroup}{Group: 6}
\newcommand{\reportdate}{November 13, 2026}
\newcommand{\pinnedcommit}{\texttt{4fbd6b8e9}}% the WezTerm snapshot the course's Understand project is built from

\newcommand{\authorentry}[2]{\begin{tabular}[t]{@{}c@{}}#1\\{\small\texttt{#2}}\end{tabular}}

% Empty frame that holds a figure's place until the diagram is added.
% Replace \figureslot{} with \includegraphics[width=\linewidth]{<file>}.
\newcommand{\figureslot}[1][2.5in]{\framebox[0.9\linewidth]{\rule{0pt}{#1}}}

\hypersetup{pdftitle={\reporttitle}, pdfauthor={\reportgroup}}

\begin{document}

% ---------- outline notes: planning pages, not part of the report ----------
\ifoutline
  \begingroup\raggedright
  % Outline notes: planning pages printed before the title page while \outlinetrue is set.
% They are not part of the report.
\section*{Outline Notes}

These pages, the italic line under each heading and every bulleted list are outline material. Each bullet is a requirement the section must meet. Write your prose below the list in your own section file and delete the list once the section is done. Setting \texttt{\textbackslash outlinefalse} in \texttt{main.tex} hides all of it at once.

\subsection*{Who Writes What}

Section files are numbered as in the group's work breakdown. The report itself numbers its sections from the Introduction, so the numbers below differ from the file numbers.

\begin{center}
\small
\begin{tabularx}{\linewidth}{@{}>{\raggedright\arraybackslash}X>{\raggedright\arraybackslash}p{0.18\linewidth}>{\raggedright\arraybackslash}p{0.14\linewidth}>{\raggedright\arraybackslash}p{0.31\linewidth}@{}}
  \toprule
  \textbf{Section in the report} & \textbf{Owner} & \textbf{Target} & \textbf{File in \texttt{sections/} (\texttt{.tex})} \\
  \midrule
  Abstract & Sydney Eichenberg & 0.25 page & \texttt{01-abstract} \\
  \ref{sec:introduction} Introduction and Project Overview & Sydney Eichenberg & 0.5 page & \texttt{02-introduction} \\
  \ref{sec:data-dictionary} Data Dictionary and Naming Conventions & Everyone; Sydney checks & 0.5 page & \texttt{03-data-dictionary} \\
  \ref{sec:conceptual} Conceptual Architecture & Jeremy McCrindle & 1.25 pages + box-and-arrow diagram & \texttt{04-conceptual-architecture} \\
  \ref{sec:derivation} Architecture Derivation Process & Jeremy McCrindle & 1 page & \texttt{05-derivation-process} \\
  \ref{sec:concrete} Concrete Architecture & Aayan A. Kader & 2.25 pages + redrawn diagrams & \texttt{06-concrete-architecture} \\
  \ref{sec:reflexion} High-Level Reflexion Analysis & Jasmine Qin & 2 pages & \texttt{07-high-level-reflexion} \\
  \ref{sec:subsystem} Second-Level Reflexion Analysis & Nigel Groen & 2.75 pages & \texttt{08-second-level-reflexion} \\
  \ref{sec:use-cases} Use Cases and Sequence Diagrams & D'Artagnan Passailaigue & 2.5 pages & \texttt{09-use-cases} \\
  \ref{sec:conclusions} Conclusions and Lessons Learned & Derek Goodman & 0.75 page & \texttt{11-conclusions-lessons-learned} \\
  \ref{sec:references} References & Everyone; Sydney checks & 0.5--1 page & \texttt{12-references} \\
  Appendix~\ref{sec:ai-report} AI Teammate Interaction and Impact & Derek Goodman & exactly 2 pages, last & \texttt{10-ai-collaboration} \\
  \bottomrule
\end{tabularx}
\end{center}

The breakdown lists the AI section as item 10, before the Conclusions and References. The course requires each report to end with its 2-page AI collaboration report, as A1 did, so it is printed last as Appendix~\ref{sec:ai-report} and keeps its file number.

\subsection*{Schedule (Toronto Time)}

\begin{center}
\small
\begin{tabularx}{\linewidth}{@{}>{\raggedright\arraybackslash}p{0.30\linewidth}X@{}}
  \toprule
  \textbf{Date} & \textbf{Milestone} \\
  \midrule
  Sat Nov 7 & First draft of every section \\
  Mon Nov 9 & Review and edits done \\
  Wed Nov 11 & Presentation ready \\
  Fri Nov 13, 11:59pm & Report PDF and slides submitted on onQ; the slides contain the URL of the recorded video, which shows both presenters' faces \\
  Sat Nov 14, 11:59pm & Peer review of every member, self included, on onQ \\
  Tue Nov 17 to Thu Nov 19, 11:59pm & Report and slides posted on the group website, not before Nov 17 \\
  Wed Nov 25, 11:59pm & Two anonymous 1-page reviews of other groups' videos (3 strengths, 3 improvements with suggestions), emailed to our TA \\
  \bottomrule
\end{tabularx}
\end{center}

\subsection*{Highest-Priority Quantitative Requirements}

\begin{center}
\small
\begin{tabular}{@{}ll@{}}
  \toprule
  \textbf{Requirement} & \textbf{Target} \\
  \midrule
  High-level conceptual architecture & about 10 components \\
  Detailed subsystem architecture & at least 4 sub-components \\
  Use cases & 2 \\
  Sequence diagrams & 2 \\
  Sequence diagram interactions & almost all labelled with concrete method calls \\
  Reflexion analyses & 2 (high-level and second-level) \\
  Git evidence & supports every divergence explained \\
  \bottomrule
\end{tabular}
\end{center}

The items most easily overlooked are the two reflexion analyses, git evidence for divergences, concrete method calls in the sequence diagrams, and a transparent account of the AI teammate's contribution.

\subsection*{Where Each Rubric Criterion Is Answered}

\begin{center}
\small
\begin{tabularx}{\linewidth}{@{}>{\raggedright\arraybackslash}X>{\raggedright\arraybackslash}p{0.24\linewidth}@{}}
  \toprule
  \textbf{Rubric criterion} & \textbf{Section} \\
  \midrule
  \multicolumn{2}{@{}l}{\textit{Quality of architecture}} \\
  High-level concrete architecture as a box-and-arrow diagram, at the right level, each component and dependency discussed and linked to the reference architecture & \ref{sec:concrete} \\
  Architecture derivation process explained & \ref{sec:derivation} \\
  Inner architecture of one subsystem and its interactions, conceptual and concrete, at least 4 sub-components & \ref{sec:subsystem} \\
  Two use cases in sequence diagrams that match the concrete architecture & \ref{sec:use-cases} \\
  Method calls in the two use cases refer to concrete methods & \ref{sec:use-cases} \\
  Reflexion analysis for the high-level architecture, based on facts from git & \ref{sec:reflexion} \\
  Reflexion analysis for the chosen second-level architecture, based on facts from git & \ref{sec:subsystem} \\
  Conceptual architecture (about 10 components) clearly referenced, with rationale for changes since A1 & \ref{sec:conceptual} \\
  Lessons learned, adapted to the project & \ref{sec:conclusions} \\
  \midrule
  \multicolumn{2}{@{}l}{\textit{Quality of report presentation}} \\
  Consistency (box names, legends, etc.) & Whole report; \ref{sec:data-dictionary} \\
  Screenshots of Understand with readable fonts, edges and arrows, possibly redrawn & \ref{sec:concrete}, \ref{sec:subsystem} \\
  Good subsystem layouts, easy to compare conceptual and concrete & \ref{sec:conceptual}--\ref{sec:subsystem} \\
  Abstract gives a clear overview & Abstract \\
  Good organization, within page limits & \ref{sec:introduction}; page targets \\
  Diagrams and tables explained in the text; diagrams have legends & Every figure and table \\
  References listed together and cited from the text & \ref{sec:references} \\
  Interaction with the AI teammate (selection, tasks, protocol, prompting) & Appendix~\ref{sec:ai-report} \\
  Quality and impact of the AI teammate (quality control, reflection, quantitative impact) & Appendix~\ref{sec:ai-report} \\
  \bottomrule
\end{tabularx}
\end{center}

\subsection*{Rules for Everyone}

\begin{itemize}[itemsep=1pt]
  \item Write only in your own section file. Add terms to the Data Dictionary and entries to the References as you use them, since both are shared.
  \item Use the component names in the Data Dictionary everywhere: diagrams, sequence diagrams and text.
  \item Cite sources as \texttt{\textbackslash cite\{key\}}, with the entry added to \texttt{12-references.tex}. Copy any A1 entry you reuse from \texttt{../a1-report/sections/references.tex}.
  \item Put figures in \texttt{figures/}, with file names starting with your section's file number (for example \texttt{06-concrete-understand.png}), so names do not collide.
  \item Give every figure a legend, and refer to every figure and table in the text (``Figure~\textbackslash ref\{fig:...\} shows \ldots'').
  \item Take every git fact (commit hash, author, date) from a clone checked out at \pinnedcommit, so it matches the snapshot Understand analyses.
  \item Log every AI interaction (tool and version, prompt, what was kept, how it was checked) and send the logs to Derek for the AI appendix. The logs are submitted with the report.
\end{itemize}

\subsection*{Format and Open Questions}

\begin{itemize}[itemsep=1pt]
  \item The format follows our A1 report: 10pt, slightly tightened line spacing, Latin Modern, black only. The A1 guidelines asked for 12pt, single-spaced, at most 17 pages including the 2-page AI report. The page targets above total about 16.5 pages.
  \item Not in hand yet: the detailed A2 statement that the onQ A2 page refers to, which may set the page limit and require sections the rubric does not list; and the reference architecture document on onQ, which Sections~\ref{sec:conceptual} and~\ref{sec:concrete} must link components to.
\end{itemize}

\subsection*{Understand Setup}

\begin{enumerate}[itemsep=1pt]
  \item Register for an Understand licence at \url{https://licensing.scitools.com/student} with a \texttt{@queensu.ca} address.
  \item Run \texttt{git clone git@github.com:wezterm/wezterm.git}, then \texttt{git checkout 4fbd6b8e9} inside the clone.
  \item Download \texttt{WezTerminal.und} (a zip) from the A2 page on onQ and unzip it inside the clone.
  \item In Understand, choose File $>$ Open $>$ Project, select the \texttt{.und} file, and wait for the first analysis to finish.
\end{enumerate}

  \clearpage
  \endgroup
\fi

% ---------- title, authors, abstract ----------
\begin{center}
  {\LARGE\bfseries \reporttitle\par}
  \vspace{8pt}
  {\large \reportcourse\par}
  {\large \reportassignment\par}
  \vspace{6pt}
  \reportgroup\quad\textbullet\quad\reportdate\par
  \vspace{14pt}
  \begin{tabular}{@{}ccc@{}}
    \authorentry{Aayan A. Kader}{23bbm12@queensu.ca} &
    \authorentry{D'Artagnan Passailaigue}{23vtw1@queensu.ca} &
    \authorentry{Jeremy McCrindle}{23bvj8@queensu.ca} \\ \noalign{\vspace{8pt}}
    \authorentry{Nigel Groen}{nigel.groen@queensu.ca} &
    \authorentry{Derek Goodman}{23kwk5@queensu.ca} &
    \authorentry{Sydney Eichenberg}{22bl71@queensu.ca} \\ \noalign{\vspace{8pt}}
    & \authorentry{Jasmine Qin}{jasmine.qin@queensu.ca} &
  \end{tabular}
\end{center}
\vspace{6pt}

\section*{Abstract}
\label{sec:abstract}

\sectionnote{Owner: Sydney Eichenberg. Target: 0.25 page, on the title page. Rubric: Abstract gives a clear overview of the report.}

\begin{requirements}
  \item ``Great'' on this criterion means a clear, easy-to-follow abstract that matches the report's contents.
  \item State what WezTerm is in one sentence, and that this report recovers its concrete architecture at commit \pinnedcommit{} with SciTools Understand and compares it with our conceptual architecture.
  \item Summarize the conceptual architecture and the main changes since A1.
  \item Name the top-level concrete subsystems and the most important divergences, with the cause git history gives for each.
  \item Name the subsystem analysed at the second level and its main finding.
  \item Name the two use cases traced through the concrete architecture.
  \item Write it last, once every other section is final, so every claim matches the body. Leave out citations, figure numbers and details a reader needs the body to understand.
\end{requirements}


% ---------- report body ----------
\section{Introduction and Project Overview}
\label{sec:introduction}

\sectionnote{Owner: Sydney Eichenberg. Target: 0.5 page. Rubric: Good report organization break-down, within page limits.}

\begin{requirements}
  \item ``Great'' on this criterion means a clear structure and logical flow, with no repetition, within the page limit.
  \item A reader of the introduction alone should know what the report found and where each part is.
\end{requirements}

\subsection{Purpose and Scope}
\label{sec:purpose}

\begin{requirements}
  \item State the goal of the report: recover WezTerm's concrete architecture, then identify and explain where it diverges from the conceptual architecture.
  \item State exactly what was analysed: the WezTerm repository at commit \pinnedcommit{} (August~31, 2026) and the course-provided Understand project built from it, both cited.
  \item State what the analysis builds on: our A1 conceptual architecture, cited, as revised in Section~\ref{sec:conceptual}.
  \item State the scope: the whole system at the top level, and one subsystem (Section~\ref{sec:subsystem}) at the second level.
\end{requirements}

\subsection{Report Organization}
\label{sec:organization}

\begin{requirements}
  \item One sentence per section, in the order they appear, ending with the AI appendix.
\end{requirements}

\subsection{Summary of Findings}
\label{sec:summary-findings}

\begin{requirements}
  \item The three or four most important findings: how closely the concrete architecture matches the conceptual one, the most significant divergences and their causes, and the main result for the second-level subsystem.
  \item Keep these consistent with the Abstract and the Conclusions.
\end{requirements}

\section{Data Dictionary and Naming Conventions}
\label{sec:data-dictionary}

\sectionnote{Owner: everyone adds entries; Sydney Eichenberg checks it. Target: 0.5 page. Rubric: Consistency (box names, legends, etc.).}

\begin{requirements}
  \item ``Great'' on this criterion means everything is consistent throughout: box names, legends and terms match in every diagram and in the text.
  \item Both tables start from A1's. Add each new term or convention when you first use it, and update an A1 entry whose meaning changed in Section~\ref{sec:conceptual}.
  \item Final check: every box name in every diagram appears in the Data Dictionary with the same spelling, and every abbreviation used in the report appears in the Naming Conventions.
\end{requirements}

\subsection{Data Dictionary}
\label{sec:dictionary}

\begin{requirements}
  \item Add the A2 analysis terms: concrete architecture, reflexion model, convergence, divergence, absence, crate.
  \item Add any new conceptual components from Section~\ref{sec:conceptual} and the sub-components from Section~\ref{sec:subsystem}.
  \item Add the code elements the report names, with their type (for example \texttt{Mux}, \texttt{Domain}, \texttt{LocalPane}, \texttt{ClientDomain}, \texttt{RemoteSshDomain}).
\end{requirements}

\begin{table}[H]
  \centering
  \caption{Data dictionary.}
  \label{tab:data-dictionary}
  \small
  \begin{tabularx}{\linewidth}{@{}>{\raggedright\arraybackslash}p{0.24\linewidth}>{\raggedright\arraybackslash}p{0.24\linewidth}X@{}}
    \toprule
    \textbf{Term} & \textbf{Type} & \textbf{Definition} \\
    \midrule
    GUI Front End & WezTerm subsystem & Displays terminal content and receives user input through graphical windows. \\
    Multiplexer & WezTerm subsystem & Coordinates windows, tabs, panes, and domains, routing requests between the GUI Front End and the relevant local or remote session. \\
    Mux Client & WezTerm subsystem & Manages connections between the local application and multiplexing domains, sending input and receiving pane updates; for remote SSH domains, it communicates with the remote Mux Server. \\
    Mux Server & WezTerm subsystem & Headless WezTerm instance that manages remote session state, including windows, tabs, panes, and their programs. \\
    Configuration Engine & WezTerm subsystem & Evaluates Lua configuration and supplies runtime settings, key bindings, themes, and event-handling support to the GUI Front End and Multiplexer. \\
    SSH Client & WezTerm subsystem & Establishes the authenticated, encrypted SSH connection used to reach a remote Mux Server. \\
    SSH Server & External system & The remote \texttt{sshd} service that accepts SSH connections and authenticates users. \\
    Terminal Emulation Core & WezTerm subsystem & Interprets terminal output, including escape sequences, and updates a pane's logical screen state. \\
    PTY & Operating-system interface & Pseudo-terminal interface that carries input and output bytes between WezTerm and a shell or other program. \\
    Shell & External program & Command-line program running in a PTY that receives commands and produces terminal output. \\
    Domain & Session/execution context & Determines where a pane's program runs; panes belong to a domain, which may be local or remote. \\
    SSH multiplexing domain & Remote domain type & SSH domain that connects to a remote Mux Server, allowing its panes to keep running independently of the local GUI connection. \\
    Pane & User-interface element & A terminal area within a tab. \\
    Tab & User-interface element & A group of one or more panes within a WezTerm window. \\
    Window & User-interface element & A top-level graphical window that can contain tabs and panes. \\
    predicted echo & Latency feature & A local prediction of typed input, displayed to reduce perceived delay during slower remote sessions. \\
    \bottomrule
  \end{tabularx}
\end{table}

\subsection{Naming Conventions}
\label{sec:naming-conventions}

\begin{requirements}
  \item Replace A1's rule for sequence messages (the last row) with the A2 rule: messages are the concrete method calls, written \texttt{Type::method()}.
  \item Add a row for code identifiers in monospace exactly as in the source: crate names as in \texttt{Cargo.toml} (\texttt{wezterm-client}), file paths (\texttt{mux/src/ssh.rs}) and methods (\texttt{Domain::attach()}).
  \item Add a row for the reflexion diagrams: the line style that marks a convergence, a divergence and an absence.
  \item Add a row saying how a reader matches a conceptual component to its code, pointing to the mapping table in Section~\ref{sec:derivation}.
  \item Add the new component names to the ``Component names'' row.
\end{requirements}

To keep the report and diagrams consistent, we use the following naming conventions:

\begin{table}[!htbp]
  \centering
  \caption{Naming conventions and abbreviations.}
  \label{tab:naming-conventions}
  \begin{tabularx}{\linewidth}{@{}p{0.30\linewidth}X@{}}
    \toprule
    \textbf{Convention} & \textbf{Meaning} \\
    \midrule
    WezTerm & Always capitalize the product name as \textit{WezTerm}. \\
    Component names & Use title case for \textit{GUI Front End}, \textit{Multiplexer}, \textit{Mux Client}, \textit{Mux Server}, \textit{SSH Client}, and \textit{Terminal Emulation Core}. \\
    Mux & Use \textit{Mux} only in the names \textit{Mux Client} and \textit{Mux Server}; call the coordinating subsystem the \textit{Multiplexer}. \\
    Abbreviations & Capitalize \textit{GUI} (Graphical User Interface), \textit{SSH} (Secure Shell), \textit{PTY} (pseudo-terminal), and \textit{OS} (Operating System). \\
    External participants & In sequence diagrams, mark non-WezTerm participants with the UML stereotype \texttt{\guillemotleft{}external\guillemotright{}}, as for the Shell and SSH Server. In prose, identify them as external when relevant. \\
    Generic interface terms & Keep \textit{window}, \textit{tab}, \textit{pane}, \textit{domain}, and \textit{session} lowercase except at the start of a sentence. \\
    Dependency arrows & Solid arrows in the dependency diagram point from a component to a component it depends on; arrow labels state the dependency or flow. \\
    Sequence messages & Use concise, action-oriented lowercase labels, such as \textit{attach domain}, \textit{push changed lines}, and \textit{repaint changed lines}. Preserve capitalization for API or event identifiers such as \texttt{DetachDomain} and \texttt{PaneOutput}. \\
    \bottomrule
  \end{tabularx}
\end{table}

\section{Conceptual Architecture}
\label{sec:conceptual}

\sectionnote{Owner: Jeremy McCrindle. Target: 1.25 pages plus the box-and-arrow diagram. Rubric: Conceptual architecture breakdown is clearly referenced, including rationale for updating the architecture (compared to A1); Good subsystem layouts.}

\subsection{Starting Point: The Conceptual Architecture Report}
\label{sec:a1-summary}

Our conceptual architecture report~\cite{a1-report} described WezTerm through eight subsystems, identified by tracing SSH multiplexing use cases. The GUI Front End, Multiplexer, Mux Client and SSH Client ran on the local host, and the Mux Server, Terminal Emulation Core and PTY ran on the remote host. A Configuration Engine was described in the text but not drawn. The report said little about fonts, windowing or scripting, which every session uses, so this report extends the architecture to the whole system.

\subsection{Changes Since the Conceptual Architecture Report}
\label{sec:a1-changes}

Each change is based on WezTerm's documentation, the conceptual architecture reports of other groups, or inconsistencies in our own report. None is based on the source code, because differences from the code are the divergences analysed in Sections~\ref{sec:reflexion} and~\ref{sec:subsystem}. From group 22~\cite{g22-a1} we adopted font handling as a separate component and configuration as a component that scripts can call, and from group 13~\cite{g13-a1}, domains, which make local and remote panes behave alike. Table~\ref{tab:components} lists the change to each component; none was removed.

\begin{figure}[!bt]
  \centering
  \begin{minipage}{0.85\linewidth}\diagram{conceptual-architecture.tex}\end{minipage}
  \caption{Revised conceptual architecture of WezTerm. An arrow points from a component to the component it depends on, and dashed arrows are configuration reads. Below the client-server line, the server runs its own Multiplexer, Terminal Emulation Core and PTY, which are drawn once. For readability, the figure omits User $\rightarrow$ Mux Client (\texttt{wezterm cli}), GUI Front End $\rightarrow$ Mux Server and three configuration reads, all of which are listed in Table~\ref{tab:dependencies}.}
  \label{fig:conceptual}
\end{figure}

The main dependency change concerns the PTY. Our earlier report showed the Terminal Emulation Core depending on the PTY~\cite{a1-report}. The Multiplexer now depends on it instead, since the domain that creates a pane also creates its PTY~\cite{what-is-a-terminal,multiplexing}, and all three other groups model it the same way~\cite{g7-a1,g13-a1,g22-a1}. The Terminal Emulation Core only reads and writes bytes, so the earlier arrow showed a flow of data, not a dependency. For the same reason, the Mux Server now depends on the Multiplexer instead of the Terminal Emulation Core.

The revision also fixes two inconsistencies in our earlier report. The Configuration Engine was described but not drawn, and it is now drawn. The local-pane use case placed a Mux Server and Terminal Emulation Core on the local host, while the figure showed both only on the remote host. The revised figure has no host boundaries, since the same components serve local panes and run inside a remote Mux Server.

\subsection{Revised Conceptual Architecture}
\label{sec:revised-conceptual}

Figure~\ref{fig:conceptual} shows the 11 components and five external entities, with a local command following the centre column and the remote path on the right. The architecture keeps the two styles of our earlier report and adds a third~\cite{arch-styles}. In the client-server style, the Mux Client and SSH Client are clients of the Mux Server and SSH Server, and the Mux Protocol defines how the Mux Client and Mux Server communicate~\cite{multiplexing}. In the implicit invocation style, the GUI Front End and Multiplexer announce events, such as new output in a pane, to listeners they do not know~\cite{event-registration}, so the figure has no arrow from the Multiplexer to the GUI Front End. The new interpreter style covers the Configuration and Lua Engine, an embedded Lua 5.4 interpreter that runs the user's \texttt{wezterm.lua}~\cite{lua-reference}.

Table~\ref{tab:components} compares each component with the web-browser reference architecture~\cite{refarch-browser}, the closest published one. Like a browser, a terminal parses marked-up content into a model, displays it and can be extended by scripts. WezTerm differs in reaching its content sources through the Multiplexer, the counterpart of the browser engine, and in having a server side. Table~\ref{tab:dependencies} explains each dependency. Dependencies 14 and 19 are inferred from the documentation.

\begin{table}[!htbp]
  \centering
  \caption{Components of the revised conceptual architecture, with each component's change from the conceptual architecture report.}
  \label{tab:components}
  \footnotesize
  \setlength{\tabcolsep}{4pt}
  \begin{tabularx}{\linewidth}{@{}>{\raggedright\arraybackslash}p{0.215\linewidth}>{\raggedright\arraybackslash}X>{\raggedright\arraybackslash}p{0.23\linewidth}>{\raggedright\arraybackslash}p{0.15\linewidth}@{}}
    \toprule
    \textbf{Component} & \textbf{Role} & \textbf{Change} & \textbf{Browser counterpart} \\
    \midrule
    GUI Front End & Handles input and draws panes & Narrowed & User Interface \\
    Window and Platform Layer & Provides native windows and events & New; was inside the GUI & Display Backend \\
    Font Subsystem & Finds, shapes and rasterises fonts & New; split from the GUI & Display Backend \\
    Multiplexer & Holds the sessions and creates panes & Now owns local panes & Browser Engine \\
    Terminal Emulation Core & Parses each pane's bytes into a cell grid & Runs with every Multiplexer & Rendering Engine \\
    PTY & Runs each program in a pseudo-terminal & Now also used locally & Networking \\
    Configuration and Lua Engine & Runs \texttt{wezterm.lua} and its handlers & Renamed and now drawn & JavaScript Interpreter \\
    Mux Client & Displays a Mux Server's panes locally & Widened beyond SSH & Networking \\
    Mux Protocol & Defines the client-server messages & New; formerly implicit & Networking \\
    Mux Server & Serves a Multiplexer's panes to clients & Serves its own Multiplexer & Web server \\
    SSH Client & Built-in SSH client (libssh or libssh2) & Adds plain SSH domains & Networking \\
    \bottomrule
  \end{tabularx}
\end{table}

\begin{table}[!htbp]
  \centering
  \caption{Dependencies of the revised conceptual architecture. Dependencies 15, 16 and 19--23 are configuration reads.}
  \label{tab:dependencies}
  \footnotesize
  \begin{tabularx}{\linewidth}{@{}l>{\raggedright\arraybackslash}p{0.53\linewidth}>{\raggedright\arraybackslash}X@{}}
    \toprule
    \textbf{\#} & \textbf{Dependency} & \textbf{Reason} \\
    \midrule
    1 & GUI Front End $\rightarrow$ Multiplexer & The GUI displays and controls the session \\
    2, 3 & GUI Front End $\rightarrow$ Window and Platform Layer, Font Subsystem & The GUI draws text in native windows \\
    4 & Multiplexer $\rightarrow$ Terminal Emulation Core & Pane output updates the terminal model \\
    5 & Multiplexer $\rightarrow$ PTY & The local domain runs programs in PTYs \\
    6, 7 & Multiplexer $\rightarrow$ Mux Client, SSH Client & Remote and SSH domains supply panes \\
    8 & Mux Client $\rightarrow$ Multiplexer & The client registers remote panes \\
    9 & Mux Client $\rightarrow$ Mux Server & The server owns the remote sessions \\
    10, 11 & Mux Client, Mux Server $\rightarrow$ Mux Protocol & Both use the same message format \\
    12 & Mux Server $\rightarrow$ Multiplexer & The server applies client requests \\
    13 & Mux Client $\rightarrow$ SSH Client & SSH multiplexing domains use SSH \\
    14 & GUI Front End $\rightarrow$ Mux Server & The GUI serves \texttt{wezterm cli} \\
    15, 16 & GUI Front End, Multiplexer $\rightarrow$ Configuration and Lua Engine & Settings and Lua event handlers \\
    17, 18 & Configuration and Lua Engine $\rightarrow$ Multiplexer, GUI Front End & Scripts call their Lua APIs \\
    19--23 & Font Subsystem, Terminal Emulation Core, Mux Client, Mux Server, SSH Client $\rightarrow$ Configuration and Lua Engine & Each reads its settings, such as the colour palette \\
    \bottomrule
  \end{tabularx}
\end{table}

\section{Architecture Derivation Process}
\label{sec:derivation}
% \depends{...}: a small inline marker for text that waits on another section. Delete each once resolved.
\providecommand{\depends}[1]{{\small\itshape\guillemotleft depends on #1\guillemotright}}

\sectionnote{Owner: Jeremy McCrindle. Target: 1 page, including the mapping table. Rubric: Architecture derivation process explained.}

Following the process taught in the course~\cite{reflexion-models} and used by Bowman et al.\ for Linux~\cite{linux-concrete}, we mapped every source file to a component of Section~\ref{sec:conceptual}, used Understand to aggregate the dependencies between components, and compared the result with the conceptual architecture.

\subsection{Tools and Setup}
\label{sec:setup}

We analysed WezTerm at commit \pinnedcommit{}~\cite{wezterm-repo} with SciTools Understand 8.0 (build 1266)~\cite{understand}. The course-provided Understand project could not complete its re-analysis on our machine, as it required more than 16~GB of memory, mainly for the vendored C libraries under \texttt{deps/}. We therefore analysed the 442 Rust files of the same commit in batches. The only errors, from one macro in the Terminal Emulation Core, affect no dependency between components. All git history was taken from a full clone at the same commit.

\subsection{Mapping Source Code to Conceptual Components}
\label{sec:mapping}

We created one Understand architecture node for each component, using the same names, and two further nodes so that every file is assigned: Shared Utilities, for generic libraries that no component owns, and Excluded. The first pass assigned each of the 64 Cargo crates to the component whose responsibilities its description and main types implement. For example, \texttt{mux} went to the Multiplexer and \texttt{codec} to the Mux Protocol. The second pass moved four files whose main type belongs to a different component, including the domain set-up in \texttt{wezterm-mux-server-impl}, which went to the Multiplexer, and the \texttt{wezterm.gui} Lua module, which went to the Configuration and Lua Engine. A file was not moved because it used another component's types, since that use is a dependency. Generated code, tests, build scripts, vendored libraries and the shared launcher were excluded. Two later iterations corrected the analysis and left the mapping unchanged (Section~\ref{sec:extracting}). Table~\ref{tab:mapping} shows the final mapping.

\begin{table}[!htbp]
  \centering
  \caption{Mapping of source code to components (lines of hand-written Rust).}
  \label{tab:mapping}
  \footnotesize
  \begin{tabularx}{\linewidth}{@{}>{\raggedright\arraybackslash}p{0.27\linewidth}>{\raggedright\arraybackslash}Xr@{}}
    \toprule
    \textbf{Component} & \textbf{Crates and folders} & \textbf{Lines} \\
    \midrule
    GUI Front End & \texttt{wezterm-gui}, \texttt{lfucache}, \texttt{frecency} & 37,830 \\
    Window and Platform Layer & \texttt{window} & 22,464 \\
    Font Subsystem & \texttt{wezterm-font} & 10,266 \\
    Multiplexer & \texttt{mux}, \texttt{bintree}, \texttt{procinfo} & 13,571 \\
    Terminal Emulation Core & \texttt{term}, \texttt{wezterm-escape-parser}, \texttt{vtparse} and 3 smaller crates & 28,966 \\
    PTY & \texttt{pty} & 2,477 \\
    Configuration and Lua Engine & \texttt{config}, \texttt{luahelper}, \texttt{lua-api-crates/} (15 crates) & 14,196 \\
    Mux Client & \texttt{wezterm-client}, \texttt{ratelim}, \texttt{wezterm/src/cli/} & 6,304 \\
    Mux Protocol & \texttt{codec} & 1,265 \\
    Mux Server & \texttt{wezterm-mux-server}, \texttt{wezterm-mux-server-impl} & 2,200 \\
    SSH Client & \texttt{wezterm-ssh} & 5,934 \\
    Shared Utilities & \texttt{termwiz}, \texttt{promise}, \texttt{filedescriptor} and 13 smaller crates & 20,803 \\
    \bottomrule
  \end{tabularx}
\end{table}

\subsection{Extracting the Concrete Dependencies}
\label{sec:extracting}

Understand aggregates file-level references, such as calls, type uses, imports and trait implementations, into dependencies between architecture nodes. We exported these through its Python API, giving 60 dependencies supported by 19,636 references. Two errors in Understand's Rust analysis assigned calls to \texttt{anyhow}'s \texttt{.context()}, and unqualified calls such as \texttt{eq}, to unrelated WezTerm code. Removing these 2,750 references removed seven false dependencies. The remaining dependencies match the crate graph from the \texttt{Cargo.toml} files in 38 cases, and each of the six differences is explained by a moved file or by a declared dependency that is never used. Section~\ref{sec:concrete} redraws the graph using the box positions of Figure~\ref{fig:conceptual}. \depends{Section~\ref{sec:concrete}}

\subsection{Reflexion Analysis Procedure}
\label{sec:reflexion-procedure}

Sections~\ref{sec:reflexion} and~\ref{sec:subsystem} classify each dependency as a convergence, which appears in both architectures, a divergence, which appears only in the concrete architecture, or an absence, which appears only in the conceptual architecture~\cite{reflexion-models}. For each divergence, we determined which code causes it, who introduced it, when, and why~\cite[slides 15--20]{reflexion-models}. We selected a representative reference in Understand, used \texttt{git blame} and \texttt{git log -S} (with \texttt{--follow}) to find the commit that introduced it, and read that commit and any linked pull request~\cite[slides 39--43]{reflexion-models}. For each absence, \texttt{git log -S} showed whether the dependency had ever existed. We recorded the reference, commit, author, date and stated rationale, and assessed whether the rationale justifies the dependency. Sections~\ref{sec:reflexion} and~\ref{sec:subsystem} explain \depends{Sections~\ref{sec:reflexion} and~\ref{sec:subsystem}: number} divergences, selected by \depends{Sections~\ref{sec:reflexion} and~\ref{sec:subsystem}: selection rule}.

\section{Concrete Architecture}
\label{sec:concrete}

\sectionnote{Owner: Aayan A. Kader. Target: 2.25 pages plus redrawn diagrams. Rubric: Detailed descriptions of the high-level concrete architecture as box-and-arrow diagram; Screenshots of Understand with readable fonts, edges and arrows, possibly redrawn; Good subsystem layouts.}

\begin{requirements}
  \item ``Great'' on these criteria means the architecture is at the right level, every component and dependency is discussed in terms of what it represents and why it is there, with a link to the reference architecture document on onQ, and the diagrams are carefully laid out and easy to compare with the conceptual architecture.
  \item The breakdown lists ``+ redrawn use case/sequence diagrams'' for this section, while the sequence diagrams and the use case diagram belong to Section~\ref{sec:use-cases} (D'Artagnan). Agree with D'Artagnan who redraws which diagram. This outline assumes this section redraws the Understand dependency graph.
\end{requirements}

\subsection{Box-and-Arrow Diagram}
\label{sec:concrete-diagram}

\begin{requirements}
  \item A screenshot of Understand's architecture dependency graph for the mapping in Section~\ref{sec:mapping}, exported at high resolution and cropped, with readable fonts and visible edges and arrows.
  \item A redrawn version that keeps every box where it sits in the conceptual figure (Section~\ref{sec:revised-conceptual}), with no overlapping boxes and as few edge crossings as possible. Optionally label each edge with Understand's dependency count.
  \item A legend for each figure, and text that walks the reader through it.
  \item Stay at the subsystem level: no files, classes or functions in this diagram.
\end{requirements}

\subsection{Subsystems}
\label{sec:concrete-subsystems}

\begin{requirements}
  \item For each concrete subsystem: what it represents (its crates and folders, from Section~\ref{sec:mapping}), its role in the running system, why it is there, its conceptual counterpart, and its counterpart in the reference architecture document on onQ, cited.
  \item Name the key code elements that carry each role, for example \texttt{Mux}, \texttt{Domain} and \texttt{LocalPane} for the Multiplexer.
  \item Name any subsystem with no conceptual counterpart, such as the shared protocol crate \texttt{codec} or a group of utility crates, and say how the analysis treats it.
\end{requirements}

\subsection{Dependencies}
\label{sec:concrete-dependencies}

\begin{requirements}
  \item One table with a row per edge in the figure: from, to, what the dependency represents with one example from the code (file and function), and why it exists.
  \item Mark each row as a convergence or divergence, and point to Section~\ref{sec:reflexion} for the explanation.
\end{requirements}

\section{High-Level Reflexion Analysis}
\label{sec:reflexion}

\sectionnote{Owner: Jasmine Qin. Target: 2 pages. Rubric: Detailed reflexion analysis for high-level architecture (for a given divergence, not all calls or messages have to be explained).}

\begin{requirements}
  \item ``Great'' on this criterion means a clear explanation of the reflexion analysis and a detailed discussion of the divergences between the conceptual and concrete architecture, based on facts obtained from git. The rubric marks down divergences that are not discussed in detail or that do not rely on actual facts from git.
  \item Follow the procedure in Section~\ref{sec:reflexion-procedure} and refer to it rather than repeating it.
\end{requirements}

\subsection{Reflexion Model}
\label{sec:reflexion-model}

\begin{requirements}
  \item One figure in the conceptual layout with every dependency marked as a convergence, divergence or absence by line style (black only, so it reads in greyscale), and a legend.
  \item A summary of the counts: how many convergences, divergences and absences.
\end{requirements}

\subsection{Divergences}
\label{sec:divergences}

\begin{requirements}
  \item One paragraph or subsection per divergence, each giving: the dependency (component A to component B); one representative example from the code (file, function, \texttt{Type::method()}); when and by whom it was introduced (commit hash, author, date, and commit message or pull request, from git); why, according to that history; and whether it is a reasonable design decision or architectural erosion.
  \item Explain one or two representative calls per divergence, not every call.
  \item Hypotheses to check first. Each crate's \texttt{Cargo.toml} at \pinnedcommit{} declares these crate-level dependencies. They are not yet confirmed in Understand, and Understand's graph is what the report must show:
  \begin{itemize}[itemsep=1pt,topsep=1pt]
    \item Reversed arrow: \texttt{wezterm-client} (Mux Client) depends on \texttt{mux} (Multiplexer), and \texttt{mux} does not depend on \texttt{wezterm-client}; A1 drew the Multiplexer depending on the Mux Client. \texttt{ClientDomain} in \texttt{wezterm-client/src/domain.rs} implements the \texttt{mux} crate's \texttt{Domain} trait.
    \item Multiplexer to SSH Client: \texttt{mux} depends on \texttt{wezterm-ssh}, through \texttt{RemoteSshDomain} in \texttt{mux/src/ssh.rs}.
    \item Multiplexer to PTY and to Terminal Emulation Core: \texttt{mux} depends on \texttt{portable-pty} (folder \texttt{pty}) and \texttt{wezterm-term} (folder \texttt{term}). \texttt{LocalPane} in \texttt{mux/src/localpane.rs} holds both the pane's \texttt{Terminal} and its \texttt{MasterPty}.
    \item GUI Front End to Terminal Emulation Core, SSH Client, Mux Client and Mux Server: \texttt{wezterm-gui} depends on \texttt{wezterm-term}, \texttt{wezterm-ssh}, \texttt{wezterm-client} and \texttt{wezterm-mux-server-impl}, not only on \texttt{mux} and \texttt{config}.
    \item Mux Server to Mux Client: \texttt{wezterm-mux-server-impl} depends on \texttt{wezterm-client}.
    \item Configuration Engine to Terminal Emulation Core and SSH Client: \texttt{config} depends on \texttt{wezterm-term} and \texttt{wezterm-ssh}.
  \end{itemize}
\end{requirements}

\subsection{Absences}
\label{sec:absences}

\begin{requirements}
  \item One short paragraph per absence: the conceptual dependency with no counterpart in the code, and why (the dependency runs through another component, or the conceptual model was wrong).
  \item Hypothesis to check first: \texttt{wezterm-term} (Terminal Emulation Core) does not depend on \texttt{portable-pty} (PTY), so A1's arrow from the Terminal Emulation Core to the PTY probably has no direct counterpart.
\end{requirements}

\section{Second-Level Reflexion Analysis}
\label{sec:subsystem}

\sectionnote{Owner: Nigel Groen. Target: 2.75 pages. Rubric: Detailed explanations of the inner architecture of one subsystem and its interactions (both conceptual and concrete); Detailed reflexion analysis for chosen 2nd-level architecture.}

\begin{requirements}
  \item ``Great'' on the inner-architecture criterion means both the conceptual and the concrete architecture of the subsystem are clearly decomposed, with at least 4 sub-components. The rubric marks down a conceptual architecture given ``without much explanation (possibly identified after the fact)''.
  \item ``Great'' on the reflexion criterion is the same as for Section~\ref{sec:reflexion}: a clear explanation and detailed discussion of each divergence, based on facts obtained from git.
  \item Choose one subsystem and justify the choice in the text. A strong candidate is the Multiplexer (crate \texttt{mux}): both use cases pass through it, it joins the local and remote sides, and its source already separates into more than four parts: the \texttt{Mux} registry and its notifications (\texttt{lib.rs}), the \texttt{Domain} abstraction and \texttt{LocalDomain} (\texttt{domain.rs}), the SSH domain (\texttt{ssh.rs}), the local pane runtime (\texttt{localpane.rs}), the window, tab and pane model (\texttt{window.rs}, \texttt{tab.rs}, \texttt{pane.rs}) and tmux integration (\texttt{tmux*.rs}).
  \item Name the chosen subsystem in the section heading.
\end{requirements}

\subsection{Conceptual Architecture of the Subsystem}
\label{sec:subsystem-conceptual}

\begin{requirements}
  \item Derive it before mapping files in Understand, from the documentation (for the Multiplexer, the multiplexing pages and the \texttt{wezterm.mux} Lua API) and the subsystem's description in Section~\ref{sec:conceptual}. Say so in the text.
  \item At least 4 sub-components, each with its responsibility, its role and the reason it exists.
  \item One figure with a legend: the sub-components, their dependencies, and the top-level components they interact with, drawn outside the subsystem's frame.
\end{requirements}

\subsection{Concrete Architecture of the Subsystem}
\label{sec:subsystem-concrete}

\begin{requirements}
  \item One table assigning the subsystem's files to the sub-components of Section~\ref{sec:subsystem-conceptual}, with the rule used for any file that was hard to place.
  \item An Understand screenshot and a redrawn diagram in the same layout as the conceptual figure, both readable and with legends.
  \item For each concrete sub-component: what it contains, what it is for, and its dependencies, including those that leave the subsystem.
  \item Connect each concrete sub-component to its conceptual counterpart.
\end{requirements}

\subsection{Reflexion Analysis of the Subsystem}
\label{sec:subsystem-reflexion}

\begin{requirements}
  \item The procedure of Section~\ref{sec:reflexion-procedure} and the level of detail of Section~\ref{sec:reflexion}, at this level: a reflexion figure or table, then each divergence with a representative code example, the git facts (commit hash, author, date, commit message or pull request) and the rationale, then the absences.
\end{requirements}

\section{Use Cases and Sequence Diagrams}
\label{sec:use-cases}

\sectionnote{Owner: D'Artagnan Passailaigue. Target: 2.5 pages, including one use case diagram and two sequence diagrams. Rubric: Two use cases are presented in sequence diagrams and match the concrete architecture; The method calls in the two use cases refer to concrete methods.}

\begin{requirements}
  \item ``Great'' on these criteria means two use cases whose sequence diagrams match the components and dependencies of the concrete box-and-arrow diagram, with almost every edge labelled with a concrete method call.
  \item Participants are boxes of the concrete architecture (Section~\ref{sec:concrete-diagram}), with identical names. External participants keep the \guillemotleft external\guillemotright{} stereotype from A1.
  \item Every message between two participants follows a dependency in the concrete diagram. Check each sequence diagram against the concrete diagram before finalising it.
  \item Label almost every edge with the method actually called at \pinnedcommit, written \texttt{Type::method()}, for example \texttt{Domain::attach()}, \texttt{Pane::key\_down()}, \texttt{Mux::add\_pane()} or \texttt{Mux::notify()}. Trace them with Understand's call graphs. Abstract labels from A1 such as ``attach domain'' or ``push changed lines'' no longer meet the rubric.
  \item Each diagram has a legend and is explained in the text with numbered steps, as in A1.
  \item Reuse the two A1 use cases so the change from the conceptual flow to the concrete one is visible, and say for each where the concrete flow differs from A1's diagram.
\end{requirements}

\subsection{Use Case Diagram}
\label{sec:use-case-diagram}

\begin{requirements}
  \item One UML use case diagram showing the actors (the user, and any external systems such as the SSH Server) and the two use cases traced below, with a legend.
  \item A1's use case diagram for SSH multiplexing (\texttt{a1-report/figures/ssh-mux-use-case.png}) can be the starting point.
\end{requirements}

\subsection{Use Case 1: Executing a Command in a Local Pane}
\label{sec:use-case-1}

\begin{requirements}
  \item The scenario: the actor, the trigger and the outcome.
  \item The participants, which are boxes of the concrete architecture.
  \item The sequence diagram, its legend, and a numbered walkthrough naming each method call.
  \item How the concrete flow differs from A1's Use Case~1, which routed a local pane through a ``local Mux Server''.
\end{requirements}

\subsection{Use Case 2: Sending Keyboard Input to a Remote Pane over SSH}
\label{sec:use-case-2}

\begin{requirements}
  \item The scenario: the actor, the trigger and the outcome.
  \item The participants, which are boxes of the concrete architecture, including the SSH Client and the Mux Server.
  \item The sequence diagram, its legend, and a numbered walkthrough naming each method call, including the call that sends the key press across the connection and the call that delivers the changed lines back.
  \item How the concrete flow differs from A1's diagrams for Use Case~2. If attaching to the remote host is the stronger example in the code, trace that instead and say why.
\end{requirements}

\section{Conclusions and Lessons Learned}
\label{sec:conclusions}

\sectionnote{Owner: Derek Goodman. Target: 0.75 page. Rubric: Lessons learned are discussed.}

\subsection{Conclusions}
\label{sec:conclusions-findings}

\begin{requirements}
  \item A summary of the key findings: how closely the concrete architecture follows the conceptual one, the most significant divergences and what git history says caused them, and the result for the second-level subsystem.
  \item What the findings mean for WezTerm's maintainability and further evolution.
  \item Proposals for future work, for example how the findings shape the enhancement proposal in A3.
  \item Keep these consistent with the Abstract and the Introduction's summary of findings.
\end{requirements}

\subsection{Lessons Learned}
\label{sec:lessons-learned}

\begin{requirements}
  \item ``Great'' on this criterion means interesting, detailed lessons adapted to the project's setting. The rubric marks down lessons that are abstract or ``copied from course notes without adapting''.
  \item Relate each lesson directly to this analysis, for example: mapping a Rust workspace of many crates onto conceptual components in Understand, tracing the commit behind a dependency in git, keeping diagram layouts aligned across the report, and what the divergences taught us about our A1 model.
  \item What the group would do differently and what it wishes it had known at the start. Gather one or two from every member.
\end{requirements}


% ---------- references ----------
\section{References}\label{sec:references}
\begingroup
\renewcommand{\section}[2]{}% the numbered heading above replaces the list's own heading
\small
% References for the report, cited in the text with \cite{key} and listed in IEEE style, as in A1.
% Everyone adds the entries they cite; Sydney checks the list. Copy any A1 entry you reuse from
% ../a1-report/sections/references.tex.
\begin{requirements}
  \item Owner: everyone adds entries; Sydney Eichenberg checks the list. Target: 0.5--1 page.
  \item Rubric: References inserted where necessary. ``Great'' means references are listed together and cited from within the text.
  \item IEEE style, numbered in order of first citation, with links, as in A1.
  \item Include: the A1 report; the WezTerm repository at \pinnedcommit; SciTools Understand; the reference architecture document on onQ; the course material on concrete architecture and reflexion models; every documentation page cited; and every commit, pull request or issue used as evidence for a divergence.
  \item Every entry must be cited in the text; remove any that are not.
\end{requirements}

\begin{thebibliography}{99}

\bibitem{a1-report}
A. A. Kader, D. Passailaigue, J. McCrindle, N. Groen, D. Goodman, S. Eichenberg, and J. Qin,
``WezTerm Conceptual Architecture,'' CISC 322/326 Assignment 1, Group 6, Queen's University, Oct. 9, 2026.
[Online]. Available: \url{https://drive.google.com/file/d/1vbQ63_ROQ0MDfKPQxsrhfRRKZknO6o8L/view?usp=sharing}.

\bibitem{wezterm-repo}
W. Furlong \emph{et al.}, ``wezterm,'' \emph{GitHub Repository}, commit 4fbd6b8e90e2326b8e25c589768f98bd71ddd047, Aug. 31, 2026.
[Online]. Available: \url{https://github.com/wezterm/wezterm/tree/4fbd6b8e90e2326b8e25c589768f98bd71ddd047}.
[Accessed: Nov. 7, 2026].

\bibitem{understand}
SciTools, ``Understand,'' version 8.0 (Build 1266).
[Online]. Available: \url{https://scitools.com/}.

\bibitem{g22-a1}
Group 22 (The Zucks), ``Conceptual Architecture of WezTerm,'' CISC 322/326 Assignment 1, Queen's University, Oct. 9, 2026.
[Online]. Available: \url{https://cisc322.netlify.app/assets/a1-report.pdf}.
[Accessed: Nov. 7, 2026].

\bibitem{g13-a1}
Group 13 (Fritz the Coders), ``CISC 322 A1 Report: Conceptual Architecture of WezTerm,'' CISC 322/326 Assignment 1, Queen's University, Oct. 2026.
[Online]. Available: \url{https://cisc-326-website.gabe-dimovski.workers.dev/CISC%20322%20A1%20Report.pdf}.
[Accessed: Nov. 7, 2026].

\bibitem{multiplexing}
W. Furlong, ``Multiplexing,'' WezTerm Documentation.
[Online]. Available: \url{https://wezterm.org/multiplexing.html}.
[Accessed: Nov. 7, 2026].

\bibitem{what-is-a-terminal}
W. Furlong, ``What is a Terminal?,'' WezTerm Documentation.
[Online]. Available: \url{https://wezterm.org/what-is-a-terminal.html}.
[Accessed: Nov. 7, 2026].

\bibitem{lua-reference}
W. Furlong, ``Lua Reference,'' WezTerm Documentation.
[Online]. Available: \url{https://wezterm.org/config/lua/general.html}.
[Accessed: Nov. 7, 2026].

\bibitem{g7-a1}
Group 7 (WezLearn), ``WezTerm Conceptual Architecture,'' CISC 322/326 Assignment 1, Queen's University, Oct. 2026.
[Online]. Available: \url{https://drive.google.com/file/d/1q7SXwHpHlMfkHrS6tZGqRjOP4EOsvnDJ/view}.
[Accessed: Nov. 7, 2026].

\bibitem{arch-styles}
A. E. Hassan, modified by B. Adams, ``Architecture Styles,'' Module 03 course slides, CISC 322/326 Software Architecture, Queen's University, Fall 2026.

\bibitem{event-registration}
W. Furlong, ``wezterm.on,'' WezTerm Documentation.
[Online]. Available: \url{https://wezterm.org/config/lua/wezterm/on.html}.
[Accessed: Nov. 7, 2026].

\bibitem{refarch-browser}
A. Grosskurth and M. W. Godfrey, ``A reference architecture for web browsers,'' in \emph{Proc. 21st IEEE Int. Conf. Software Maintenance (ICSM)}, 2005, pp. 661--664, doi: 10.1109/ICSM.2005.13.

\bibitem{reflexion-models}
A. E. Hassan, modified by B. Adams, ``Reflexion Models,'' Module 06 course slides, CISC 322/326 Software Architecture, Queen's University, Fall 2026.

\bibitem{linux-concrete}
I. T. Bowman, R. C. Holt, and N. V. Brewster, ``Linux as a case study: Its extracted software architecture,'' in \emph{Proc. 21st Int. Conf. Software Engineering (ICSE)}, 1999, pp. 555--563, doi: 10.1145/302405.302691.

\end{thebibliography}
\endgroup

% ---------- AI collaboration report: the last two pages ----------
\clearpage
\appendix
\section{AI Teammate Interaction and Impact}
\label{sec:ai-report}

\sectionnote{Owner: Derek Goodman; every member supplies AI-use logs. Target: exactly the last 2 pages. Rubric: Interaction with AI teammate; Quality and impact of AI teammate.}

\begin{requirements}
  \item ``Great'' on the interaction criterion means AI model selection, task assignment and the AI interaction protocol, including the prompting strategy, are explained well enough that another group could reproduce them. ``Great'' on the impact criterion means a clear and reasonable quality-control procedure, an insightful reflection on human--AI team dynamics, and transparency about the AI's quantitative impact on the report.
  \item Build it from the interaction logs every member keeps from the first prompt, and submit the logs with the report.
  \item Show that AI use followed a deliberate plan rather than being ad hoc.
\end{requirements}

\subsection{AI Member Profile and Selection Process}
\label{sec:ai-profile}

\begin{requirements}
  \item The name and version of every AI tool used.
  \item Why each was the right choice for A2's tasks, compared with at least one alternative the group considered. In A1 we wrote that we did not compare tools; the rubric asks for model selection to be explained.
\end{requirements}

\subsection{Tasks Assigned to the AI Teammate}
\label{sec:ai-tasks}

\begin{requirements}
  \item An itemised list of the tasks given to the AI, who assigned each, and why it suited an AI. Likely A2 tasks: summarising the git history behind a divergence, explaining Rust traits and modules found in Understand, proposing mapping rules for files that are hard to place, and drafting or revising text.
\end{requirements}

\subsection{Interaction Protocol and Prompting Strategy}
\label{sec:ai-protocol}

\begin{requirements}
  \item Who prompted and the workflow, described well enough that another group could reproduce it.
  \item At least one complex or critical prompt, quoted, and how it was refined (context given, persona, output format).
\end{requirements}

\subsection{Validation and Quality Control Procedures}
\label{sec:ai-validation}

\begin{requirements}
  \item The specific checks applied to AI output, for example: every dependency the AI claimed was confirmed in Understand; every commit hash, author and date was confirmed in a clone at \pinnedcommit; every method name in a sequence diagram was confirmed in the source.
  \item Who carried out each check, and one example of an AI error the checks caught.
\end{requirements}

\subsection{Quantitative Contribution to the Final Deliverable}
\label{sec:ai-contribution}

\begin{requirements}
  \item A percentage estimate of the AI's contribution to the submitted report, with a table breaking it down by area, and how the estimate was made (for example, from the logs).
  \item The estimate must hold up against the submitted logs; the rubric marks down an impact that is ``missing or highly inaccurate''.
\end{requirements}

\subsection{Reflection on Human--AI Team Dynamics}
\label{sec:ai-reflection}

\begin{requirements}
  \item The AI's strengths and limitations on this deliverable: whether it saved time or created new work, how it influenced decisions, and any disagreements it caused.
  \item What changed in how we worked with the AI since A1, and what we will keep for A3.
\end{requirements}


\end{document}
